\chapter{Managing Data Frames: Import, Missing Values, Filtering, Sorting, and Export}
\label{chap:managing-data-frames}

\chapterguide
  {You should be able to create, inspect, and index data frames.}
  {install and load the file/data-management packages shown in the handout; identify the working directory; import Excel and CSV files; count and locate missing values; remove rows containing missing values; rename columns; filter and search records; sort rows; and export results.}

\section{Packages used in Lab 7a}

The handout introduces several packages as the workflow expands:

\begin{dstable}
\begin{tabularx}{\linewidth}{@{}p{3.0cm}X@{}}
\toprule
\rowcolor{DSPaleBlue!92}
Package & Supplied use \\
\midrule
\texttt{readxl} & Import the Excel file \texttt{uforeport.xls}. \\
\texttt{dplyr} & Filter, pipe, and arrange data frames. \\
\texttt{xlsx} & Export an Excel file. \\
\bottomrule
\end{tabularx}
\end{dstable}

Installation and loading follow the pattern introduced in Lab 1:

\begin{lstlisting}[style=dsr]
install.packages("readxl")
library(readxl)

any(grepl("readxl", installed.packages()))
\end{lstlisting}

The \texttt{any(grepl(...))} line is the handout's explicit package-installation check.

\section{Working directory and file placement}

Before importing a file, Lab 7a asks you to inspect or set the working directory:

\begin{lstlisting}[style=dsr]
getwd()

# Example only; your path may differ.
setwd("C:/Users/User/Desktop")

list.files()
\end{lstlisting}

The handout expects \texttt{uforeport.xls} and \texttt{titanic.csv} to be stored in the same working location as the R work so that the filenames can be used directly.

\begin{keyidea}
The import examples assume the filenames resolve from the current working directory. If R cannot find a file, first check \texttt{getwd()} and \texttt{list.files()} before changing the analysis code.
\end{keyidea}

\section{Importing Excel and CSV files}

The Excel file is read with \texttt{read\_excel()} after loading \texttt{readxl}:

\begin{lstlisting}[style=dsr]
uforeport <- read_excel("uforeport.xls")
View(uforeport)
\end{lstlisting}

The CSV file uses base R's \texttt{read.csv()} and does not require the \texttt{readxl} package:

\begin{lstlisting}[style=dsr]
titanic <- read.csv("titanic.csv")
View(titanic)
\end{lstlisting}

\section{Identifying missing values}

The handout uses \texttt{is.na()} in three ways.

\subsection{Count all missing values}

\begin{lstlisting}[style=dsr]
sum(is.na(uforeport))
sum(is.na(titanic))
\end{lstlisting}

\subsection{Find their positions}

\begin{lstlisting}[style=dsr]
which(is.na(uforeport))
which(is.na(titanic))
\end{lstlisting}

A small vector is included to make the behavior easier to see:

\begin{lstlisting}[style=dsr]
demo <- c(1, 2, NA, 4, NA, 6, 7)

sum(is.na(demo))
which(is.na(demo))
\end{lstlisting}

\subsection{Count missing values by column}

\begin{lstlisting}[style=dsr]
print(sapply(uforeport, function(x) sum(is.na(x))))
print(sapply(titanic, function(x) sum(is.na(x))))
\end{lstlisting}

This gives a per-column missing-value count rather than one total for the entire table.

\section{Removing rows with missing values}

The source compares data-frame dimensions before and after \texttt{na.omit()}:

\begin{lstlisting}[style=dsr]
dim(uforeport)
uforeport_cleaned <- na.omit(uforeport)
dim(uforeport_cleaned)

dim(titanic)
titanic_cleaned <- na.omit(titanic)
dim(titanic_cleaned)
\end{lstlisting}

The dimension comparison is the evidence requested by the exercise: it shows how many rows remain after records containing NA values are omitted.

\section{Column names and renaming}

The lab retrieves names in two equivalent-looking ways:

\begin{lstlisting}[style=dsr]
colnames(uforeport_cleaned)
colnames(titanic_cleaned)

names(uforeport_cleaned)
names(titanic_cleaned)
\end{lstlisting}

For the UFO report, spaces in column names are replaced by underscores:

\begin{lstlisting}[style=dsr]
names(uforeport_cleaned) <- gsub(
  " ", "_", colnames(uforeport_cleaned)
)
\end{lstlisting}

This prepares names such as \texttt{Colors Reported} for the later \texttt{Colors\_Reported} filter expressions shown in the handout.

\section{Filtering data}

After loading \texttt{dplyr}, Lab 7a filters for exact values and numerical conditions:

\begin{lstlisting}[style=dsr]
install.packages("dplyr")
library(dplyr)

print(filter(uforeport_cleaned, City == "Belton"))
print(filter(uforeport_cleaned, Colors_Reported == "GREEN"))
print(filter(titanic_cleaned, sex == "female"))
print(filter(titanic_cleaned, fare > 500))
\end{lstlisting}

\subsection{Multiple conditions}

The supplied examples show both explicit \texttt{\&} and the pipe form:

\begin{lstlisting}[style=dsr]
print(filter(
  titanic_cleaned,
  sex == "female" & fare > 500
))

# Alternative dplyr form from the handout
titanic_cleaned %>%
  filter(sex == "female", fare > 500)

ufo_green <- uforeport_cleaned %>%
  filter(Colors_Reported == "GREEN")
\end{lstlisting}

\section{Arranging rows}

The handout sorts Titanic by fare:

\begin{lstlisting}[style=dsr]
# Ascending
titanic_sortbyfare <- arrange(titanic_cleaned, fare)

# Descending
titanic_sortbyfare <- arrange(titanic_cleaned, desc(fare))
\end{lstlisting}

The same variable name is reassigned in the source, so after both lines the descending version is the one retained.

\section{Exporting analysis results}

Lab 7a exports one Excel file and one CSV file:

\begin{lstlisting}[style=dsr]
install.packages("xlsx")
library(xlsx)
write.xlsx(ufo_green, "ufo_green.xlsx")

write.csv(
  titanic_sortbyfare,
  "titanic_sortbyfare.csv"
)
\end{lstlisting}

This completes the full file workflow: import, inspect/clean, transform, then save the result.

\begin{quickcheck}
\begin{enumerate}
  \item Which function tells you the current working directory?
  \item Which import function is used for \texttt{uforeport.xls}, and which for \texttt{titanic.csv}?
  \item How do \texttt{sum(is.na(x))} and \texttt{which(is.na(x))} answer different questions?
  \item What evidence does the handout collect before and after \texttt{na.omit()}?
  \item Which \texttt{dplyr} function sorts rows, and how is descending order requested?
\end{enumerate}
\end{quickcheck}

\begin{chapterrecap}
\begin{itemize}
  \item File import depends on the working directory and filenames used by the code.
  \item Missing values can be counted globally, located, and summarized per column.
  \item \texttt{na.omit()} is the cleaning method demonstrated in the source material.
  \item \texttt{gsub()} is used to normalize column names before filtering.
  \item \texttt{dplyr::filter()} and \texttt{arrange()} implement row selection and sorting in the supplied workflow.
  \item The final transformed tables are exported to Excel and CSV.
\end{itemize}
\end{chapterrecap}
